\chapter{隋唐五代索引}
\label{app:index}

\section{事件索引：正文中的可核查入口}

索引只链接实际存在的正文锚点，不把账本中的年度核对项或未展开事件伪装成读者已可访问的叙事。

\begin{description}
  \item[\eventref{V03-581-ANNAL-CENTRAL}{隋受禅（581）}] 隋受禅与再统一的资源条件。
  \item[\eventref{ST-617-LI-YUAN-TAIYUAN}{太原起兵（617）}] 唐建国的军政与空间条件。
  \item[\eventref{ST-690-ZHOU-FOUNDATION}{武周建国（690）}] 改号、合法性与材料边界。
  \item[\eventref{ST-755-AN-LUSHAN-REBELLION}{安禄山起兵（755）}] 范阳南下与两京危机。
  \item[\eventref{LT-880-062}{黄巢入长安（880）}] 晚唐崩解与区域军府。
\end{description}

\section{来源索引}
\begin{description}
  \item[来源与史料说明] \hyperref[app:sources]{“来源与史料说明”}
\end{description}

\section{图表与地图索引}
\begin{description}
  \item[隋的再统一] \hyperref[fig:vol03:sui:sui-reunification]{“隋受禅、南北整合与交通关系示意”}
  \item[唐初边疆与丝路] \hyperref[fig:vol03:early-tang:tang-frontiers-silk-road]{“唐初边疆、丝路与供给关系示意”}
  \item[武周、开元与边地] \hyperref[fig:vol03:wu-zhou-xuanzong:wu-zhou-kaiyuan-centers-frontiers]{“武周开元的中心与边地关系示意”}
  \item[安史战争] \hyperref[fig:vol03:an-lushan-late-tang:tang-frontiers-an-shi]{“安史战争的边镇、道路与两京示意”}
  \item[五代十国] \hyperref[fig:vol03:five-dynasties-ten-kingdoms:late-tang-five-dynasties-regional-order]{“晚唐五代的区域政权关系示意”}
\end{description}
